% Template for the Journal of Computational Science and Applications (JCSA)
% Published by the University of Lahore -- ISSN 2710-3102
% Official author guidelines: https://journals.uol.edu.pk/JCSA/Author-Guideline
\documentclass[twoside,10pt]{article}
\usepackage{lipsum}  % For sample/placeholder text
\usepackage{fancyhdr}
\usepackage{geometry}
\usepackage{multicol}
\usepackage{abstract}
\usepackage{graphicx}
\usepackage{caption}
\usepackage{titlesec}
\usepackage[utf8]{inputenc} % For UTF-8 input encoding
\usepackage[T1]{fontenc}    % Better font encoding for PDF output
\usepackage{ragged2e}  % Provides the justifying command
\usepackage{amssymb}  % For additional symbols
\usepackage{parskip}  % Automatically handles paragraph spacing
\usepackage{amsmath}  % For advanced math environments
\usepackage{newtxtext}  % Times New Roman font
\usepackage{booktabs}  % For improved table appearance
\usepackage{array}     % For customizing table columns
\usepackage{hyperref} % For hyperlinks in the document
\usepackage[backend=biber,style=numeric,sorting=none]{biblatex} % For bibliography management


% The starting page number is normally assigned by the editorial office
% once the article is scheduled into a specific issue. Authors can leave
% this at 1; production will update it if required.
\setcounter{page}{1}

% Define custom column types for better alignment and spacing
\newcolumntype{L}[1]{>{\raggedright\arraybackslash}p{#1}} % Left aligned column
\newcolumntype{C}[1]{>{\centering\arraybackslash}p{#1}}   % Center aligned column
\newcolumntype{R}[1]{>{\raggedleft\arraybackslash}p{#1}}  % Right aligned column

\usepackage[font=footnotesize, labelfont=bf, justification=raggedright, format=plain]{caption}

\captionsetup[table]{
    labelsep=period,               % Period after label for tables
    justification=raggedright      % Ensure left alignment of captions
}

\captionsetup[figure]{
    labelsep=period,               % Period after label for figures
    justification=raggedright      % Ensure left alignment of captions
}

\addbibresource{sources.bib} %Import the bibliography file
% Customizing bibliography heading
\defbibheading{bibliography}[\refname]{%
  \section{\MakeUppercase{REFERENCES}}}

% Set page margins
\geometry{
    a4paper,
    left=0.8in,
    right=0.8in,
    top=1in,
    bottom=1in
}

% Increase space between columns
\setlength{\columnsep}{15pt}

% Define the centred abstract heading
\renewcommand{\abstractname}{\hspace{1.3em}\textbf{Abstract}}

\titleformat{\section}{\large\bfseries\uppercase}{\thesection.}{1em}{}

\titleformat{\subsection}{\normalfont\bfseries\flushleft}{\thesubsection.}{1em}{}

\titleformat{\subsubsection}{\normalfont\bfseries\itshape\raggedright}{\thesubsubsection.}{1em}{}

% Set font sizes for abstract
\renewcommand{\abstractnamefont}{\normalfont\fontsize{14}{20}\bfseries}
\renewcommand{\absnamepos}{flushleft}
\renewcommand{\abstracttextfont}{\normalfont\fontsize{11}{13}\selectfont}

% Keywords environment
\newcommand{\keywordsname}{Keywords}
\newenvironment{keywords}
  {\par\noindent\hspace{2em}\textbf{\keywordsname:}}
  {\par}

\hyphenpenalty=10000
\exhyphenpenalty=10000
\sloppy

% Set up header and footer styles
\pagestyle{fancy}
\fancyhf{}  % Clear default header and footer

% Odd page headers -- replace with a short running form of your title
\fancyhead[C]{\small Running Title of the Paper}

% Even page headers -- replace with the first author's surname and year
\fancyhead[CE]{\small Author Surname et al. (Year)}

% Footer: Page numbers only (on all pages except the first one)
\fancyfoot[C]{\thepage}

% First page custom header/footer
\fancypagestyle{firstpage}{
  \fancyhf{}  % Clear header and footer
  \fancyhead[LO]{\small Volume X, Issue Y, Year}  % Left header for first page
  \fancyhead[RO]{\small J. Comput. Sci. Appl.}  % Right header for first page
  \fancyfoot[L]{\footnotesize Received: Month Year; Revised: Month Year; Accepted: Month Year\\\textsuperscript{*}Corresponding Author: Full Name \textless{}email@example.com\textgreater{}}
   \fancyfoot[R]{\footnotesize \copyright\ Year by the Author (s). Published by JCSA,\\\textsuperscript{*} under the CC-BY license }
  \renewcommand{\headrulewidth}{0.4pt}  % Add header line
  \renewcommand{\footrulewidth}{0.4pt}  % Add footer line
}

\title{\fontsize{16}{20}\selectfont\textbf{Paper Title}}
\author{
    \fontsize{12}{14}\selectfont
    First Author\textsuperscript{1*} \\
    \fontsize{12}{14}\selectfont
    \textsuperscript{1}Department Name, University Name, City, Country 
}
\date{}

\begin{document}

% Change "Research Article" to the appropriate article type if needed
% (e.g. Review Article, Short Communication).
\vspace{-0.5cm}  % Adjust spacing as needed between header and "Research Article"
\raggedright%

{\fontsize{11}{13}\selectfont{Research Article}}
\vspace{-0.8cm}  % Adjust spacing between "Research Article" and title

% Now, insert the title and the author information
\maketitle Hypergraph RAG + sparce autoencoder, using quantum linear algebra (Complex / Hilbert space): to visualise concept drift

\thispagestyle{firstpage} % Apply custom header/footer for first page
\begin{abstract}
\fontsize{11}{13}\selectfont
\vspace{1em}
\justifying%
Place abstract text here

\vspace{-0.1cm}
\begin{keywords}
Place keywords here
\end{keywords}

\end{abstract}

\vspace{1cm}

\section{Introduction}
\fontsize{12}{14}\selectfont
Place text here

\section{Literature Review}
\fontsize{12}{14}\selectfont
~\cite{laine2025quantumllm} For semantic space modelling.

\subsection{Phase 0: Problem Definition}
\fontsize{12}{14}\selectfont
\subsubsection{Define the drift problem in one paragraph}
\subsubsection{Define the drift problem in one paragraph}
This project visualises concept drift in a HyperRAG system operating over simple medical datasets. 
As the system answers a query using a HyperRAG network, a sparse autoencoder identifies the concepts it relies on, and complex (quantum-inspired) 
linear algebra gives each concept a magnitude and an angle in Hilbert space; magnitude for how strongly the concept is present, 
angle for its relation to other concepts. Plotting these across queries turns the similarities and differences 
between concepts into visible geometry: as the query changes, the relative angles between concepts open and close, and that movement is the drift. 
The aim is to make the system's shifting internal concepts legible in a way a real-valued, magnitude-only view cannot.
We want to see how the system judges relations in a Hilbert space, and how those relations change over different queries.

\subsubsection{Identify target audience}
Researchers, engineers and artists.
\subsubsection{Define expected output (visualiser + baseline comparison)}
A 3D render (a projection of the high-dimensional Hilbert space) of concepts moving as queries change, 
alongside a baseline of the same concepts in real vector space for comparison.

\subsubsection{Write success criteria}
The Hilbert-space view (angle + magnitude) reveals relational structure or drift between concepts that the real-valued, 
magnitude-only baseline does not. Equivalently; degrading the Hilbert-space representation down to a real vector space visibly 
loses information, and the visualiser makes that lost information apparent.

\subsubsection{Record assumptions / risks}
Information loss and spurious concepts in both the real and Hilbert-space representations, and the risk of overfitting to a small dataset.

\subsection{Phase 1: Hypergraph RAG}
\fontsize{12}{14}\selectfont
\subsubsection{Hypergraph Store}
This section takes our initial data and stores it as a hypergraph, each entity, is a node and n-ary of facts are hyperedges. 
This data strcuture is built using hypernetx. This section has no quantum linear algebra as this si the store of our preliminary data structure.

\subsubsection{Embedding Layer}
Where entities and relations become embeddings, embed these in a Hilbert space.


\section{Results and Discussion}

Place text here

\section{Conclusion and Future Work}

Place text here

\section{Declaration}
\fontsize{11}{13}\selectfont

\textbf{Conflict of Interest: } Place text here

\textbf{Author Contribution:} Place text here

% Example for references section

\printbibliography% Use BibLaTeX for printing bibliography
\end{document}